\titular%
{Carathéodory e a axiomatización da termodinámica}
{Sebastián Táboas Pazo}
{divulgacion}
{}

\begin{refsection}
\begin{multicols}{2}

\subsection*{Introdución}
Quizais o espírito que une a todos os científicos sexa a pescuda de respostas
á realidade material que se nos impón, na física algúns acadaran o gozo na idea
dunha \textit{teoría do todo} que nos brinde explicación a todo o que acontece.
Outros, no entanto, pretenderan buscar un denominador común capaz de
fundamentar calquera teoría, mesmo \textit{unha teoría do todo}; este é un
espírito axiomatizador que se xesta dende a antigüidade xa con Aristóteles ou
Euclides e se estende até hoxe en día, así que tamén se pretendeu axiomatizar a
física. Na conferencia do Congreso internacional de Matemáticos de 1900, David
Hilbert propuxo os coñecidos \textit{23 problemas de
Hilbert}\footnote{Realmente só propuxera dez deles.}, entre eles o de
axiomatizar toda a física.

Os esforzos neste proxecto son moitos e en moitas ramas, mais agora destacamos
o traballo da axiomatización da termodinámica, quizais o menos frutífero de
todos eles. Nun primeiro momento, a persoa que aceptou este desafío foi o
matemático grego Constantin Carathéodory na súa obra
\cite{caratheodory.c_1909}


\subsection*{A axiomatización de Carathéodory}
No seu traballo de 1909, Carathéodory parte de tan só tres definicións
primitivas e dous axiomas para fundamentar toda a termodinámica que pioneiros
coma Kelvin, Carnot ou Clausius edificaran ao longo de máis de medio século.
Primeiramente, o matemático define a equivalencia de sistemas e as variábeis de
estado, como a caracterización dos cambios de estado. O primeiro axioma que
propón é unha reescritura do primeiro principio da termodinámica:

\begin{quotation}
\textit{Cada fase $\phi_i$ dun sistema S en equilibrio asóciase a unha función
$\varepsilon_i$ das cantidades $V_i$, $p_i$, $m_i$, que é proporcional ao
volume total $V_i$ da fase e denomínase enerxía interna da devandita fase. A
suma $\varepsilon = \sum\varepsilon_i$ de todas as fases chámase enerxía
interna do sistema. En calquera cambio de estado adiabático, o cambio de
enerxía incrementado polo traballo externo, $A$, é nulo; é dicir, en signos, se
se denotan os valores inicial e final da enerxía con} $\varepsilon$ \textit{e}
$\hat{\varepsilon} ~\textit{respectivamente}$:  $$\hat{\varepsilon}
–\varepsilon + A = 0$$
\hspace{1cm}\textit{Carathéodory, 1909}
\end{quotation}

Até aquí non hai nada novidoso, non obstante, debemos pensar que as expresións
matemáticas que frecuentemente manexamos en termodinámica son diferenciais de
Pfaff: $df=\sum x_i dX_i$, os cales no caso adiabático ($\dbar Q = 0$) supoñen
un problema de curvas características entre o estado inicial e o final. Nesta
clase de procesos, estas 1–formas verifican sempre o segundo lema de Schwarz,
polo que se trata de diferenciais exactos que podemos integrar; agora ben, en
procesos non adiabáticos isto non é así. O matemático grego é tamén autor dun
teorema que leva o seu nome: este asegúranos que para toda 1–forma non exacta
existe un factor integrante que a converte en integrábel. Así, afirmou que para
$\dbar Q = dU – \dbar W$ existe un factor integrante que é o inverso da
temperatura absoluta ($1/T$) que converte a $\dbar Q$ en exacta, este novo
diferencial é a entropía, $dS$. Desta forma, semella sólido enunciar o seguinte
axioma:

\begin{quotation}
<<En calquera entorno dun estado inicial arbitrario hai estados que son
inaccesíbeis mediante cambios de estado adiabáticos.>>

\hspace{1cm}\textit{Carathéodory, 1909}
\end{quotation}

\subsection*{Críticas}
Estes axiomas semellan constituír un análogo aos principios clásicos da
imposibilidade dos móbiles perpetuos, no caso de seren formulados sen ningunha
clase de erro. A crítica a este traballo é bastante ampla e deu un salto
exponencial a mediados do século pasado co auxe da termodinámica irreversíbel
na que se evidenciaban os problemas destes enunciados, cando menos
insuficientes, pois Carathéodory sempre restrinxírase a un marco de procesos
cuasiestáticos.

Unha das primeiras críticas que aparece vén de man de Max Planck, quen sinalara
a súa falta de contacto coa evidencia experimental nunha rama de fundamento
plenamente fenomenolóxica.

Este último sentenciou:
\begin{quote}
\textit{…ninguén até agora tratou de alcanzar soamente mediante procesos
adiabáticos todo punto no entorno de calquera estado de equilibrio, e así
comprobar se de veras son inaccesíbeis, […], este axioma non nos brinda o menor
indicio que nos permita diferenciar os estados accesíbeis dos inaccesíbeis.}

\cite{planck.m_1926}
\end{quote}

Outra crítica de entre moitas radica, aínda que se evidencie a existencia dun
factor integrante vinculado á escala absoluta Kelvin, en que este non comporta
unha función da clase das temperaturas termodinámicas, a pesar de ser este erro
rectificábel.

Finalmente, este traballo de 1909 condenouse a unha posición de mera
curiosidade académica e, polo tanto, acabou por ocupar un lugar secundario nos
manuais de termodinámica. Porén, o espírito da formalización termodinámica
pervive como imperecedoiro.

\printbibliography
\end{multicols}
\end{refsection}
